\chapter{Question on Notice on the Norwegian and European Models}
\label{ch:supp:qon}

\section{The question}
\label{sec:supp:qon:q}

\begin{quote}
\itshape
I was curious about the Norwegian model, in particular, and that data processing infrastructure that meets European standards. If you could elaborate on that a little, because what you were referring to just then with crypto, are there examples of what we were talking about in the session earlier around a kind of data supply chain, essentially, where you put in rules throughout the use of data, whether it's storage or collection or even analysis, and how that relates to the regulation of data centres in Europe?
\end{quote}

This response proceeds on two limbs. The first examines Norway's treatment of cryptocurrency mining; the second sets out the broader European standards and the data supply chain they govern. The unifying point is that Europe now regulates the data centre not as a bare energy user but as one node in a wider regulated infrastructure spanning physical siting, cybersecurity, privacy, data portability and environmental performance.

\section{The Norwegian model}
\label{sec:supp:qon:norway}

Norway distinguishes data centres that contribute to national digital infrastructure from those that merely consume large quantities of power, cryptocurrency mining being the paradigm of the latter. Under its 2025 data centre strategy, a facility exceeding 500 kilowatts must register with the government and disclose its physical location, the services it provides, specified clients, its subscribed electrical capacity, and the power it estimates is consumed in cryptocurrency mining \autocite{norwayDataCentreStrategy2025}. The policy intent is restrictive. The Government has resolved to investigate \enquote{a temporary and limited prohibition on the establishment of data centres that mainly engage in energy-intensive cryptocurrency mining}, on the view that such mining is resource intensive while generating few jobs and little income \autocite{norwayDataCentreStrategy2025}. That intent has since hardened. Following a private member's bill, in March 2026 the Norwegian Parliament asked the Government to return with a formal legislative proposal for a ban on establishing data centres used for cryptocurrency mining \autocite{schjodtNorwayCryptoResolutions2026}.

\section{European standards and the data supply chain}
\label{sec:supp:qon:europe}

The European Union takes a similar but broader approach. Under the Energy Efficiency Directive it now requires monitoring and reporting of the energy performance and water footprint of significant data centres, and a data centre energy efficiency package is in preparation that will assess the reported data and introduce a rating scheme for European facilities \autocite{ecEnergyPerformanceDataCentres2024}. This is the same measurement-and-disclosure foundation recommended in \Cref{ch:supp:pue}, applied at supranational scale.

On the data supply chain the question raises, two instruments divide the work. The General Data Protection Regulation governs how personal data may be handled, while the Data Governance Act governs the intermediaries that handle protected public-sector data. Europe pursues this regulation from a position of structural dependence, since \enquote{just three US-based companies account for 65\% of the EU cloud services market} \autocite{marcelinCloudAIDevelopment2025}, alongside constrained access to water, land and capital for new market entrants.

The proposed response is a Cloud and AI Development Act, expected to take the form of a directly applicable European regulation \autocite{aliEUCloudAIAct2025}. It would define what qualifies as highly secure cloud infrastructure for critical uses such as defence, public administration and critical infrastructure, and it aims to \enquote{triple the EU's data centre capacity within the next five to seven years} by streamlining permitting and easing access to land, energy and capital, with a fast-track approval process for projects assessed to carry high social impact \autocite{marcelinCloudAIDevelopment2025}. For the Committee, the Norwegian and European models converge on the sequencing this submission urges for New South Wales, in which registration, measurement and disclosure are the precondition for any allocative or prohibitory lever.
